\begin{trapbox}
	\begin{description}[style=nextline, leftmargin=2.8em, labelsep=0.5em, itemsep=0.35em]
		\item[\textbf{\textcolor{CTrap}{易错点一（误把连心线中点当成弦中点）：}}] 看到 $AB\perp O_1O_2$，就认为 $M$ 一定是 $O_1O_2$ 的中点。
		\item[\textbf{\textcolor{CTrap}{正确理解（半径不同则通常偏移）：}}] $M$ 是\textbf{公共弦 $AB$ 的中点}，不是必然的两圆心线段中点。事实上 $t=(d^2+r_1^2-r_2^2)/(2d)$；只有 $r_1=r_2$ 时才有 $t=d/2$。
		\item[\textbf{\textcolor{CTrap}{错因分析（混淆两条线段）：}}] 把“连心线是弦的垂直平分线”误读成“公共弦经过连心线的中点”。
	\end{description}

	\medskip
	\begin{description}[style=nextline, leftmargin=2.8em, labelsep=0.5em, itemsep=0.35em]
		\item[\textbf{\textcolor{CTrap}{易错点二（默认中点夹在两圆心之间）：}}] 不加判断地设 $O_2M=d-O_1M$，并把所有距离都看作非负。
		\item[\textbf{\textcolor{CTrap}{正确理解（使用有向位置）：}}] 定义 $t$ 为沿 $O_1\to O_2$ 的有向投影坐标，两个勾股式写成 $r_1^2=t^2+h^2$、$r_2^2=(d-t)^2+h^2$。可能出现 $t<0$ 或 $t>d$；几何长度 $O_2M=|d-t|$。
		\item[\textbf{\textcolor{CTrap}{错因分析（把直线误当线段）：}}] 两圆确有两个交点时，公共弦中点也可能在两圆心连线段的外部。\textbf{例：}$r_1=8,r_2=5,d=4$，有 $t=55/8>4$。
	\end{description}

	\medskip
	\begin{description}[style=nextline, leftmargin=2.8em, labelsep=0.5em, itemsep=0.35em]
		\item[\textbf{\textcolor{CTrap}{易错点三（机械使用斜率负倒数）：}}] 当连心线水平或竖直时，仍然硬套 $k_{AB}=-1/k_{O_1O_2}$。
		\item[\textbf{\textcolor{CTrap}{正确理解（改用法向式）：}}] 连心线水平，则公共弦为竖直线 $x=x_M$；连心线竖直，则公共弦为水平线 $y=y_M$。统一使用 $\Delta a(x-x_M)+\Delta b(y-y_M)=0$ 最稳妥。
		\item[\textbf{\textcolor{CTrap}{错因分析（忽略斜率不存在）：}}] 竖直线没有斜率，不能用“负倒数”在所有情况下直接计算。
	\end{description}

	\medskip
	\begin{description}[style=nextline, leftmargin=2.8em, labelsep=0.5em, itemsep=0.35em]
		\item[\textbf{\textcolor{CTrap}{易错点四（把整条直线当作公共点集）：}}] 把两圆方程相减得到的一次方程，误认为其上所有点都在两个圆上。
		\item[\textbf{\textcolor{CTrap}{正确理解（仅保证交点落在线上）：}}] 两圆交点 $A,B$ 满足相减后的方程，故它是过 $A,B$ 的直线；但直线上其他点通常不在任意一个圆上。
		\item[\textbf{\textcolor{CTrap}{错因分析（误用逆命题）：}}] $f_1=f_2=0\Rightarrow f_1-f_2=0$，而 $f_1-f_2=0$ 不能单独推出 $f_1=f_2=0$。
	\end{description}

	\medskip
	\begin{description}[style=nextline, leftmargin=2.8em, labelsep=0.5em, itemsep=0.35em]
		\item[\textbf{\textcolor{CTrap}{易错点五（没有检查相交条件）：}}] 在两圆相离、内含或相切时，直接把相减所得直线叫作“经过两个公共交点的公共弦直线”。
		\item[\textbf{\textcolor{CTrap}{正确理解（先确认相交于两点）：}}] 只有 $|r_1-r_2|<d<r_1+r_2$ 才存在非零长度的公共弦 $AB$。相切时只有一个公共点；不相交时没有实公共弦。圆心重合时须另作判断。
		\item[\textbf{\textcolor{CTrap}{错因分析（忘记适用范围）：}}] 相减得到一次式并不自动保证两个圆存在两个交点。
	\end{description}
\end{trapbox}
